\documentclass[11pt]{article}
\newcommand{\ListaTitulo}{Lista 10: Bernoulli, Binomial e Poisson}
% Preâmbulo compartilhado: MATF14, Listas de Exercícios 2026
% Usado por ListaNN.tex e GabaritoNN.tex via % Preâmbulo compartilhado: MATF14, Listas de Exercícios 2026
% Usado por ListaNN.tex e GabaritoNN.tex via % Preâmbulo compartilhado: MATF14, Listas de Exercícios 2026
% Usado por ListaNN.tex e GabaritoNN.tex via \input{preamble}

\usepackage[utf8]{inputenc}
\usepackage[T1]{fontenc}
\usepackage[brazil]{babel}
\usepackage{fullpage}
\usepackage{amsmath,amssymb}
\usepackage{enumitem}
\usepackage{xcolor}
\usepackage{fancyhdr}
\usepackage{titlesec}
\usepackage{listings}
\usepackage{hyperref}
\usepackage{booktabs}
\usepackage{graphicx}
\usepackage{tikz}

\definecolor{matf14azul}{HTML}{0D3B54}
\definecolor{matf14dourado}{HTML}{C9971E}

\titleformat{\section}{\normalfont\Large\bfseries\color{matf14azul}}{\thesection}{1em}{}
\titleformat{\subsection}{\normalfont\large\bfseries\color{matf14azul}}{\thesubsection}{1em}{}

\providecommand{\ListaTitulo}{Lista de Exercícios}

\setlength{\headheight}{14pt}
\pagestyle{fancy}
\fancyhf{}
\lhead{\small MATF14: Estatística Econômica I}
\rhead{\small \ListaTitulo}
\cfoot{\thepage}
\renewcommand{\headrulewidth}{0.4pt}

\lstset{
  basicstyle=\ttfamily\small,
  breaklines=true,
  frame=single,
  columns=fullflexible,
  backgroundcolor=\color{gray!8},
  commentstyle=\color{matf14dourado},
  keywordstyle=\color{matf14azul}\bfseries,
  language=R
}

\newcounter{questaocounter}
\newenvironment{questao}[1]{%
  \stepcounter{questaocounter}%
  \bigskip\noindent\textbf{Questão #1.}\ %
}{\par}

\newcommand{\cabecalho}[2]{%
  \begin{center}
    {\Large\bfseries\color{matf14azul} #1}\\[0.3em]
    {\large #2}
  \end{center}
  \vspace{0.5em}
  \noindent\rule{\linewidth}{0.6pt}
  \vspace{0.5em}
}

\newenvironment{solucao}{%
  \par\smallskip\noindent\textit{\color{matf14dourado!80!black}Solução.}\ %
}{\hfill$\blacksquare$\par}


\usepackage[utf8]{inputenc}
\usepackage[T1]{fontenc}
\usepackage[brazil]{babel}
\usepackage{fullpage}
\usepackage{amsmath,amssymb}
\usepackage{enumitem}
\usepackage{xcolor}
\usepackage{fancyhdr}
\usepackage{titlesec}
\usepackage{listings}
\usepackage{hyperref}
\usepackage{booktabs}
\usepackage{graphicx}
\usepackage{tikz}

\definecolor{matf14azul}{HTML}{0D3B54}
\definecolor{matf14dourado}{HTML}{C9971E}

\titleformat{\section}{\normalfont\Large\bfseries\color{matf14azul}}{\thesection}{1em}{}
\titleformat{\subsection}{\normalfont\large\bfseries\color{matf14azul}}{\thesubsection}{1em}{}

\providecommand{\ListaTitulo}{Lista de Exercícios}

\setlength{\headheight}{14pt}
\pagestyle{fancy}
\fancyhf{}
\lhead{\small MATF14: Estatística Econômica I}
\rhead{\small \ListaTitulo}
\cfoot{\thepage}
\renewcommand{\headrulewidth}{0.4pt}

\lstset{
  basicstyle=\ttfamily\small,
  breaklines=true,
  frame=single,
  columns=fullflexible,
  backgroundcolor=\color{gray!8},
  commentstyle=\color{matf14dourado},
  keywordstyle=\color{matf14azul}\bfseries,
  language=R
}

\newcounter{questaocounter}
\newenvironment{questao}[1]{%
  \stepcounter{questaocounter}%
  \bigskip\noindent\textbf{Questão #1.}\ %
}{\par}

\newcommand{\cabecalho}[2]{%
  \begin{center}
    {\Large\bfseries\color{matf14azul} #1}\\[0.3em]
    {\large #2}
  \end{center}
  \vspace{0.5em}
  \noindent\rule{\linewidth}{0.6pt}
  \vspace{0.5em}
}

\newenvironment{solucao}{%
  \par\smallskip\noindent\textit{\color{matf14dourado!80!black}Solução.}\ %
}{\hfill$\blacksquare$\par}


\usepackage[utf8]{inputenc}
\usepackage[T1]{fontenc}
\usepackage[brazil]{babel}
\usepackage{fullpage}
\usepackage{amsmath,amssymb}
\usepackage{enumitem}
\usepackage{xcolor}
\usepackage{fancyhdr}
\usepackage{titlesec}
\usepackage{listings}
\usepackage{hyperref}
\usepackage{booktabs}
\usepackage{graphicx}
\usepackage{tikz}

\definecolor{matf14azul}{HTML}{0D3B54}
\definecolor{matf14dourado}{HTML}{C9971E}

\titleformat{\section}{\normalfont\Large\bfseries\color{matf14azul}}{\thesection}{1em}{}
\titleformat{\subsection}{\normalfont\large\bfseries\color{matf14azul}}{\thesubsection}{1em}{}

\providecommand{\ListaTitulo}{Lista de Exercícios}

\setlength{\headheight}{14pt}
\pagestyle{fancy}
\fancyhf{}
\lhead{\small MATF14: Estatística Econômica I}
\rhead{\small \ListaTitulo}
\cfoot{\thepage}
\renewcommand{\headrulewidth}{0.4pt}

\lstset{
  basicstyle=\ttfamily\small,
  breaklines=true,
  frame=single,
  columns=fullflexible,
  backgroundcolor=\color{gray!8},
  commentstyle=\color{matf14dourado},
  keywordstyle=\color{matf14azul}\bfseries,
  language=R
}

\newcounter{questaocounter}
\newenvironment{questao}[1]{%
  \stepcounter{questaocounter}%
  \bigskip\noindent\textbf{Questão #1.}\ %
}{\par}

\newcommand{\cabecalho}[2]{%
  \begin{center}
    {\Large\bfseries\color{matf14azul} #1}\\[0.3em]
    {\large #2}
  \end{center}
  \vspace{0.5em}
  \noindent\rule{\linewidth}{0.6pt}
  \vspace{0.5em}
}

\newenvironment{solucao}{%
  \par\smallskip\noindent\textit{\color{matf14dourado!80!black}Solução.}\ %
}{\hfill$\blacksquare$\par}


\begin{document}

\cabecalho{Lista de Exercícios 10}{Distribuições de Bernoulli, Binomial e Poisson (Cap. 3, Aulas 20--21)}

\noindent \textit{Justifique todas as respostas. As resoluções são discutidas em sala e nos horários de atendimento; não são publicadas neste material.}

\begin{questao}{1} \textbf{(Binomial: carteira de investidores)}

Uma corretora observa que 30\% de seus clientes de varejo realizam pelo menos uma operação por
semana, de forma independente entre clientes. Em um grupo de 12 clientes sorteados:

\begin{enumerate}[label=(\alph*)]
  \item Qual é a probabilidade de exatamente 4 realizarem operação na semana?
  \item Qual é a probabilidade de \emph{no máximo} 2 realizarem operação?
  \item Calcule $E(X)$ e $dp(X)$, o número esperado e o desvio padrão de clientes ativos no grupo.
\end{enumerate}
\end{questao}

\begin{questao}{2} \textbf{(Binomial: controle de qualidade)}

Um lote de peças tem 4\% de defeito, independente entre peças. Um comprador aceita o lote se, em
uma amostra de 25 peças testadas, houver \textbf{no máximo 1} defeituosa.

\begin{enumerate}[label=(\alph*)]
  \item Qual a probabilidade de o lote ser aceito?
  \item Se a taxa de defeito real do fornecedor subir para 10\% (sem o comprador saber), qual
    passa a ser a probabilidade de aceitação do mesmo plano de amostragem? O plano de amostragem
    ficou mais ou menos eficaz em rejeitar lotes ruins?
\end{enumerate}
\end{questao}

\begin{questao}{3} \textbf{(Poisson: fila de atendimento)}

Uma agência bancária recebe, em média, 6 clientes por hora em um determinado horário do dia,
seguindo uma distribuição de Poisson.

\begin{enumerate}[label=(\alph*)]
  \item Qual a probabilidade de a agência receber exatamente 4 clientes em uma hora?
  \item Qual a probabilidade de receber mais de 8 clientes em uma hora?
  \item Qual a probabilidade de receber exatamente 2 clientes em \textbf{meia hora}? (Dica: a taxa média
    de uma Poisson é proporcional ao tamanho do intervalo de tempo considerado.)
\end{enumerate}
\end{questao}

\begin{questao}{4} \textbf{(Aproximação Binomial $\to$ Poisson: fraude rara)}

Uma financeira processa 5.000 transações por dia; cada uma tem probabilidade $0{,}0006$ de ser
fraudulenta, independentemente das demais.

\begin{enumerate}[label=(\alph*)]
  \item Calcule $P(X=3)$ usando o modelo Binomial exato, com $X\sim\text{Bin}(5000;\,0{,}0006)$.
  \item Calcule $P(X=3)$ usando a aproximação de Poisson, com $\lambda=np$.
  \item Compare os dois resultados. Por que a aproximação funciona bem neste caso? (Relacione com
    as condições vistas em aula: $n$ grande, $p$ pequeno.)
\end{enumerate}
\end{questao}

\begin{questao}{5} \textbf{(Dedução: esperança da Binomial a partir da soma de Bernoullis)}

Seja $X\sim\text{Bin}(n,p)$, escrita como $X=X_1+X_2+\cdots+X_n$, onde cada $X_i$ é
Bernoulli($p$) independente, com $E(X_i)=p$ e $\mathrm{Var}(X_i)=p(1-p)$.

\begin{enumerate}[label=(\alph*)]
  \item Usando a propriedade $E(X_1+\cdots+X_n) = E(X_1)+\cdots+E(X_n)$ (válida sempre, mesmo sem
    independência), mostre que $E(X)=np$.
  \item Usando a propriedade $\mathrm{Var}(X_1+\cdots+X_n)=\mathrm{Var}(X_1)+\cdots+
    \mathrm{Var}(X_n)$ (válida sob independência), mostre que $\mathrm{Var}(X)=np(1-p)$.
  \item Essa dedução mostra por que a condição de \textbf{independência} entre as tentativas é
    essencial para a fórmula $\mathrm{Var}(X)=np(1-p)$, mas não é necessária para $E(X)=np$. Onde,
    exatamente, na dedução de (a) e (b), a independência é usada (ou não)?
\end{enumerate}
\end{questao}

\bigskip
\noindent\textit{A questão a seguir não tem resposta única e não é cobrada em prova no mesmo
formato: fecha a primeira metade da Unidade 3 e é para discussão em grupo ou por escrito.}

\begin{questao}{6 (discussão)} \textbf{(Sem resposta única)}

Escolha uma variável de contagem do seu interesse (número de clientes que cancelam um serviço por
mês, número de acidentes de trabalho por trimestre em uma empresa, número de transações
fraudulentas por dia). Argumente, por escrito, se Binomial ou Poisson é o modelo mais defensável
para essa variável, explicitando quais das condições de cada modelo você julga razoavelmente
satisfeitas e quais são mais frágeis na prática.
\end{questao}


\bigskip
\section*{Exercícios complementares}

\begin{enumerate}[label=\textbf{C\arabic*.}, leftmargin=*]
\item Um terço dos adultos de certa região são alfabetizados. Dentre 5 adultos escolhidos ao acaso, qual a probabilidade de
\begin{enumerate}[label=(\alph*)]
  \item dois serem alfabetizados e três não?
  \item mais de dois serem alfabetizados?
\end{enumerate}
\item Um exame tem 10 questões de verdadeiro ou falso. Um aluno responde todas ao acaso.
\begin{enumerate}[label=(\alph*)]
  \item Qual a probabilidade de acertar exatamente 7?
  \item Qual a média e o desvio padrão do número de acertos?
\end{enumerate}
\item (Distribuição geométrica) Realizam-se tentativas independentes com probabilidade de sucesso $p$. Seja $X$ o número de tentativas até o primeiro sucesso, $X \in \{1, 2, 3, \ldots\}$.
\begin{enumerate}[label=(\alph*)]
  \item Calcule $P(X = 2)$ e $P(X = 3)$.
  \item Deduza $P(X = k)$ para $k \geq 1$.
  \item Um banco de sangue procura doadores O-Rh negativo, 9\% da população. Qual a probabilidade de o primeiro doador com esse tipo ser a sexta pessoa examinada?
\end{enumerate}
\item O número médio de carros abandonados por semana em certa autoestrada é 2,2. Pelo modelo de Poisson, qual a probabilidade de
\begin{enumerate}[label=(\alph*)]
  \item nenhum carro ser abandonado na semana seguinte?
  \item pelo menos 2 carros serem abandonados?
\end{enumerate}
\item Parafusos são defeituosos com probabilidade 2\%, independentemente. São vendidos em pacotes de 10, com devolução do dinheiro se houver dois ou mais defeituosos.
\begin{enumerate}[label=(\alph*)]
  \item Qual a probabilidade de devolver o dinheiro de um pacote?
  \item Quem compra dez pacotes independentes, com que probabilidade terá dinheiro devolvido em pelo menos um?
\end{enumerate}
\item A probabilidade de um produto ter defeito é 0,01. Calcule a probabilidade de uma amostra de 200 unidades independentes ter no máximo dois defeituosos: (a) pela Binomial; (b) pela aproximação de Poisson.
\item Sabe-se que 43\% das pessoas de uma população têm sangue tipo A. Numa amostra de 120 pessoas, seja $N$ o número com sangue tipo A. Qual a distribuição de $N$, sua média e seu desvio padrão? A aproximação de Poisson seria boa aqui? Por quê?

\end{enumerate}

\bigskip
\needspace{12\baselineskip}
\section*{Aprofundamento: modelos probabilísticos e dados reais}
\noindent\textit{Exercícios ligados à seção de aplicação macroeconômica do Capítulo 3 do livro.}

\begin{enumerate}[label=\textbf{A\arabic*.}, leftmargin=*]
\item \textbf{(A Binomial serve para recessões?)} De 1997 a 2025 houve 116 trimestres, e o PIB caiu (em relação ao trimestre anterior, com ajuste sazonal)
em 26 deles. A distribuição observada do número de quedas por ano está na Lista 09 (15, 6, 5, 2 e 1 anos com 0, 1, 2, 3 e 4 quedas).
\begin{enumerate}[label=(\alph*)]
  \item Estime $p$, a probabilidade de queda em um trimestre. Supondo $X \sim \mathrm{Binomial}(4, p)$, calcule $P(X = k)$ para $k = 0, \ldots, 4$ e
    compare com as frequências relativas observadas.
  \item Compare $E(X)$ e $\mathrm{Var}(X)$ da Binomial com a média (0,897) e a variância (1,265) observadas.
  \item Qual hipótese da Binomial falha? Explique usando o que você sabe sobre crises econômicas.
\end{enumerate}

\item \textbf{(Meses de deflação e a Poisson.)} Seja $Y$ o número de meses de um ano com IPCA negativo. De 2000 a 2025 (26 anos), houve 15 anos com
nenhum mês de deflação, 8 anos com um mês, 2 anos com dois meses e 1 ano com três meses (IBGE, via BCB-SGS, série 433).
\begin{enumerate}[label=(\alph*)]
  \item Estime $\lambda$ pela média de meses de deflação por ano e compare a variância observada com $\lambda$.
  \item Supondo $Y \sim \mathrm{Poisson}(\lambda)$, calcule $P(Y = 0)$, $P(Y = 1)$ e $P(Y \geq 2)$, e compare com as frequências observadas.
  \item Dos 15 meses de deflação, 12 ocorreram de 2017 em diante. Isso é compatível com uma taxa $\lambda$ constante ao longo dos 26 anos? Que consequência isso tem para o modelo?
\end{enumerate}

\end{enumerate}

\end{document}
